\chapter{Pianificazione della traiettoria cartesiana}
\label{chap:pianificazione_traiettoria}

La pianificazione della traiettoria ha lo scopo di definire il movimento dell'organo terminale del robot nello spazio cartesiano, a partire da un insieme di punti tridimensionali già individuati sulla maglia. Questi punti rappresentano i riferimenti geometrici che il robot deve seguire durante le operazioni di cucitura simulata lungo i polsini, il fondo della maglia e il colletto.

Dal punto di vista fisico, non è sufficiente collegare direttamente i punti mediante segmenti rettilinei. Una successione di segmenti può infatti introdurre discontinuità nella direzione del movimento e richiedere variazioni istantanee della velocità in corrispondenza dei vertici. Tali condizioni non sono compatibili con un movimento regolare, poiché comporterebbero accelerazioni elevate e sollecitazioni dinamiche non desiderabili.

La traiettoria viene quindi costruita attraverso tre operazioni principali:
\begin{enumerate}
    \item organizzazione dei punti in segmenti distinti, comprendenti le cuciture e i raccordi di collegamento;
    \item interpolazione geometrica mediante spline cubiche, per ottenere percorsi continui;
    \item assegnazione di una legge temporale al movimento, attraverso profili di velocità e accelerazione limitati.
\end{enumerate}

Il risultato è una traiettoria cartesiana descritta, per ogni istante, dalla posizione desiderata dell'organo terminale, dalla sua velocità lineare e dal suo orientamento.
Grazie alla sua versatilità, il metodo realizzato può essere applicato sia al caso simulato con Unity che al caso reale con il robot UR5, con la possibilità di adattare i parametri cinematici in funzione delle esigenze operative.

\section{Organizzazione geometrica del percorso}

La funzione \texttt{genera\_traiettoria\_diversa()} organizza il percorso in otto segmenti consecutivi. Tale suddivisione permette di distinguere le operazioni di cucitura vere e proprie dai movimenti necessari per passare da una regione della maglia alla successiva.

Indicando con $\mathbf{P}_{\mathrm{start}}$ la posizione iniziale dell'organo terminale e con $\mathcal{P}_{\mathrm{sx}}$, $\mathcal{P}_{\mathrm{fondo}}$, $\mathcal{P}_{\mathrm{dx}}$ e $\mathcal{P}_{\mathrm{coll}}$ gli insiemi di punti relativi, rispettivamente, al polsino sinistro, al fondo, al polsino destro e al colletto, la sequenza è organizzata come segue:

\begin{enumerate}
    \item avvicinamento dalla posizione iniziale al primo punto del polsino sinistro;
    \item percorrenza dei punti del polsino sinistro;
    \item raccordo dal polsino sinistro al fondo maglia;
    \item percorrenza dei punti del fondo maglia;
    \item raccordo dal fondo maglia al polsino destro;
    \item percorrenza dei punti del polsino destro;
    \item raccordo dal polsino destro al colletto;
    \item percorrenza dei punti del colletto.
\end{enumerate}

Questa suddivisione consente di assegnare a ciascun tratto parametri di movimento differenti, in funzione della sua natura geometrica e delle esigenze operative.

\subsection{Interpolazione dei punti dei polsini}

Per i polsini sono inizialmente disponibili due punti estremi, che identificano l'inizio e la fine del tratto da percorrere. Per aumentare la densità dei riferimenti geometrici, vengono inseriti due punti intermedi mediante interpolazione lineare.

Indicando con $\mathbf{P}_0$ e $\mathbf{P}_1$ gli estremi del segmento, i punti intermedi sono definiti da

\begin{equation}
    \mathbf{P}_{1/3}
    = \mathbf{P}_0 + \frac{1}{3}
    \left(\mathbf{P}_1-\mathbf{P}_0\right),
\end{equation}

\begin{equation}
    \mathbf{P}_{2/3}
    = \mathbf{P}_0 + \frac{2}{3}
    \left(\mathbf{P}_1-\mathbf{P}_0\right).
\end{equation}

Si ottiene così una sequenza di quattro punti complessivi. L'inserimento di punti intermedi non modifica la geometria del segmento rettilineo iniziale, ma fornisce più riferimenti per la successiva interpolazione.

\subsection{Uniformazione della quota dei punti - Caso UR5 simulato}

Per la gestione delle coordinate lungo l'asse \textit{z} viene utilizzata distanza fissa, corrispondente a quella tra la telecamera e il maglione, precedentemente misurata attraverso script.

\subsection{Uniformazione della quota dei punti - Caso UR5 reale}

Per ridurre le variazioni di quota tra i diversi riferimenti della maglia, la coordinata verticale dei punti dei polsini e del colletto viene uniformata utilizzando una quota di riferimento calcolata a partire dai punti del fondo.

Indicando con $z_i$ la coordinata verticale dell'$i$-esimo punto del fondo, la quota comune è

\begin{equation}
    z_{\mathrm{ref}}
    = \frac{1}{N_f}\sum_{i=1}^{N_f} z_i,
    \label{eq:quota_riferimento}
\end{equation}

dove $N_f$ è il numero di punti utilizzati per descrivere il fondo.

Per i punti appartenenti ai polsini e al colletto, la coordinata verticale viene quindi sostituita con

\begin{equation}
    z_i' = z_{\mathrm{ref}}.
\end{equation}

In questo modo, i punti di tali regioni risultano complanari rispetto alla quota scelta, mentre i punti del fondo conservano le proprie coordinate. La scelta semplifica la geometria del percorso e riduce le variazioni verticali introdotte dai punti di riferimento.

È importante precisare che questa operazione costituisce un'approssimazione geometrica: non descrive direttamente la deformazione della maglia né garantisce che la sua superficie reale sia perfettamente piana.

\section{Costruzione dei raccordi tra i segmenti}

I segmenti di cucitura sono collegati mediante tratti di transizione. Lo scopo dei raccordi è evitare che il robot debba passare direttamente da una regione della maglia a un'altra seguendo un percorso che attraversi inutilmente la superficie intermedia.

La funzione \texttt{build\_trapezoids()} costruisce questi raccordi sulla base della distanza tra l'estremo del segmento precedente e l'inizio di quello successivo.

Siano $\mathbf{P}_a$ il punto finale del segmento precedente e $\mathbf{P}_b$ il punto iniziale del segmento successivo. La distanza tra i due punti, proiettati sul piano orizzontale, è calcolata come

\begin{equation}
    d_{xy}
    = \sqrt{(x_b-x_a)^2+(y_b-y_a)^2}.
    \label{eq:distanza_xy}
\end{equation}

La geometria del raccordo dipende da tale distanza:

\begin{itemize}
    \item se $d_{xy} \geq 0.1\,\mathrm{m}$, vengono introdotti due punti intermedi;
    \item se $0.03\,\mathrm{m} \leq d_{xy} < 0.1\,\mathrm{m}$, viene introdotto un solo punto intermedio;
    \item se $d_{xy} < 0.03\,\mathrm{m}$, non vengono aggiunti punti intermedi.
\end{itemize}

Nel caso di due punti intermedi, le loro coordinate orizzontali sono ottenute dividendo lo spostamento in tre parti uguali:

\begin{equation}
    \mathbf{p}_{1,xy}
    = \mathbf{p}_{a,xy}
    + \frac{1}{3}
    \left(\mathbf{p}_{b,xy}-\mathbf{p}_{a,xy}\right),
\end{equation}

\begin{equation}
    \mathbf{p}_{2,xy}
    = \mathbf{p}_{a,xy}
    + \frac{2}{3}
    \left(\mathbf{p}_{b,xy}-\mathbf{p}_{a,xy}\right).
\end{equation}

La quota dei punti intermedi viene aumentata di una quantità prefissata:

\begin{equation}
    z_{\mathrm{up}} = z_a + h,
    \qquad h=0.1\,\mathrm{m}.
    \label{eq:quota_raccordo}
\end{equation}

Quando viene inserito un solo punto intermedio, le sue coordinate orizzontali corrispondono al punto medio tra gli estremi. La quota viene calcolata con la stessa regola.

Dal punto di vista fisico, il sollevamento introduce un percorso che si allontana dalla superficie della maglia durante il trasferimento tra due regioni distanti. Questo accorgimento può ridurre il rischio di interferenza con il tessuto durante gli spostamenti, ma non costituisce una garanzia di assenza di collisioni: tale proprietà dipende anche dalla geometria effettiva della maglia, dalla posizione del tessuto e dalle dimensioni dell'organo terminale.

I raccordi vengono pertanto considerati come segmenti autonomi, ai quali è possibile assegnare limiti cinematici differenti da quelli dei tratti di cucitura.

\section{Interpolazione mediante spline cubiche}

Una volta definiti i punti di ciascun segmento, è necessario costruire una curva continua che li attraversi. A questo scopo vengono utilizzate spline cubiche naturali, calcolate separatamente per le tre coordinate cartesiane.

Sia un segmento descritto da $M$ punti tridimensionali

\begin{equation}
    \mathbf{P}_i =
    \begin{bmatrix}
        x_i & y_i & z_i
    \end{bmatrix}^{T},
    \qquad i=0,\ldots,M-1.
\end{equation}

Per associare ai punti un parametro adimensionale crescente lungo il percorso, viene calcolata la lunghezza cumulativa della poligonale:

\begin{equation}
    \ell_0 = 0,
\end{equation}

\begin{equation}
    \ell_i = \ell_{i-1}
    + \left\|\mathbf{P}_i-\mathbf{P}_{i-1}\right\|,
    \qquad i=1,\ldots,M-1.
\end{equation}

Il parametro viene poi normalizzato rispetto alla lunghezza totale della poligonale:

\begin{equation}
    u_i = \frac{\ell_i}{\ell_{M-1}},
    \qquad u_i \in [0,1].
    \label{eq:parametro_spline}
\end{equation}

La normalizzazione permette di descrivere ciascun segmento mediante un parametro che varia da zero a uno, indipendentemente dalla sua lunghezza fisica.

Vengono quindi costruite tre funzioni spline distinte:

\begin{equation}
    x=x(u), \qquad y=y(u), \qquad z=z(u).
\end{equation}

La curva cartesiana risultante è

\begin{equation}
    \mathbf{r}(u)=
    \begin{bmatrix}
        x(u)\\
        y(u)\\
        z(u)
    \end{bmatrix}.
    \label{eq:curva_cartesiana}
\end{equation}

Ogni coordinata è interpolata mediante una funzione polinomiale cubica definita a tratti. Le spline naturali impongono che la derivata seconda sia nulla agli estremi del segmento:

\begin{equation}
    x''(0)=x''(1)=0,
\end{equation}

con condizioni analoghe per $y(u)$ e $z(u)$.

La costruzione garantisce la continuità della posizione e della derivata prima all'interno del singolo segmento, evitando gli spigoli tipici della connessione diretta tra segmenti rettilinei. La derivata seconda è inoltre continua nei nodi interni della spline.

È tuttavia necessario distinguere la regolarità dei singoli segmenti dalla regolarità dell'intera traiettoria: poiché ogni tratto viene interpolato separatamente, la continuità delle derivate tra segmenti consecutivi non è automaticamente garantita dalla sola interpolazione cubica. I raccordi geometrici e le successive leggi temporali contribuiscono a rendere il movimento più regolare, ma la continuità globale deve essere valutata sui risultati ottenuti.

\section{Lunghezza della curva e parametro di avanzamento}

Il parametro $u$ descrive la geometria della spline, ma non rappresenta direttamente la distanza percorsa. Infatti, incrementi uguali di $u$ non corrispondono necessariamente a incrementi uguali di lunghezza lungo una curva.

Per associare al percorso una coordinata spaziale fisica, si considera la lunghezza d'arco. La lunghezza totale della spline è definita da

\begin{equation}
    L = \int_0^1
    \left\|\frac{d\mathbf{r}(u)}{du}\right\|du.
    \label{eq:lunghezza_arco}
\end{equation}

Nel codice, tale integrale viene approssimato numericamente campionando la curva in un elevato numero di punti e sommando le distanze euclidee tra campioni consecutivi:

\begin{equation}
    L \simeq
    \sum_{i=1}^{N_s-1}
    \left\|\mathbf{r}(u_i)-\mathbf{r}(u_{i-1})\right\|,
\end{equation}

dove $N_s$ è il numero di campioni impiegati per l'approssimazione.

Questa grandezza, espressa in metri, rappresenta la distanza effettiva lungo la curva interpolata e costituisce il riferimento per la pianificazione della velocità.

Introducendo una coordinata di avanzamento $s(t)$, con

\begin{equation}
    0 \leq s(t) \leq L,
\end{equation}

la posizione cartesiana nel tempo può essere espressa come

\begin{equation}
    \mathbf{x}(t)=\mathbf{r}\left(u(t)\right),
\end{equation}

dove $u(t)$ è ricavato dal rapporto tra l'avanzamento lungo la curva e la lunghezza totale, mediante la relazione $u(t)=s(t)/L$.

In questo modo vengono separati due problemi distinti: la definizione della forma geometrica del percorso e la definizione della velocità con cui tale percorso viene attraversato.

\section{Pianificazione della velocità e dell'accelerazione}

La sola definizione della curva non determina quanto rapidamente il robot debba percorrerla. Per descrivere il movimento nel tempo viene quindi associato a ciascun segmento un profilo di velocità trapezoidale, o triangolare quando non è disponibile spazio sufficiente per raggiungere il valore massimo di velocità.

Il profilo è definito attraverso due parametri principali:
\begin{itemize}
    \item $v_{\max}$, velocità lineare massima lungo la curva;
    \item $a_{\max}$, modulo massimo dell'accelerazione tangenziale.
\end{itemize}

La velocità scalare lungo il percorso è indicata con

\begin{equation}
    v(t)=\frac{ds(t)}{dt}.
\end{equation}

L'accelerazione tangenziale è invece

\begin{equation}
    a_t(t)=\frac{dv(t)}{dt}.
\end{equation}

Nel caso di un profilo trapezoidale, il moto è suddiviso in tre fasi:
\begin{enumerate}
    \item accelerazione, durante la quale la velocità cresce da zero fino al valore massimo;
    \item regime, durante il quale la velocità rimane costante;
    \item decelerazione, durante la quale la velocità diminuisce fino a zero.
\end{enumerate}

Assumendo accelerazione costante durante la prima fase, la velocità cresce secondo

\begin{equation}
    v(t)=a_{\max}t,
\end{equation}

mentre lo spazio percorso è

\begin{equation}
    s(t)=\frac{1}{2}a_{\max}t^2.
\end{equation}

Il tempo necessario a raggiungere la velocità massima è

\begin{equation}
    t_a=\frac{v_{\max}}{a_{\max}},
\end{equation}

e la distanza percorsa durante l'accelerazione è

\begin{equation}
    s_a=\frac{v_{\max}^2}{2a_{\max}}.
\end{equation}

Una fase di decelerazione simmetrica richiede una distanza analoga. Se la lunghezza disponibile non è sufficiente a contenere entrambe le fasi e un tratto a velocità costante, il profilo diventa triangolare: la velocità massima nominale non viene raggiunta.

La funzione di pianificazione genera un profilo spaziale di velocità lungo la curva e lo converte successivamente in una successione di istanti temporali. Il tempo viene ottenuto integrando numericamente la relazione

\begin{equation}
    dt=\frac{ds}{v(s)}.
\end{equation}

Per campioni consecutivi, l'incremento temporale viene approssimato utilizzando la velocità media:

\begin{equation}
    \Delta t_i \simeq
    \frac{s_i-s_{i-1}}
    {\left(v_i+v_{i-1}\right)/2},
    \label{eq:incremento_temporale}
\end{equation}

quando la velocità media è sufficientemente diversa da zero.

Il tempo cumulativo è quindi

\begin{equation}
    t_i=\sum_{j=1}^{i}\Delta t_j.
\end{equation}

La relazione tra tempo e avanzamento viene poi interpolata per ottenere campioni temporali uniformi. Questo passaggio consente di descrivere la traiettoria attraverso una sequenza ordinata di posizioni e velocità.

\subsection{Parametri cinematici dei segmenti}

Nel codice vengono utilizzati limiti differenti in funzione del tipo di tratto. Per i segmenti ordinari di cucitura e per i raccordi non trattati come casi speciali, i valori di riferimento sono

\begin{equation}
    v_{\max}=0.06\,\mathrm{m/s},
    \qquad
    a_{\max}=0.005\,\mathrm{m/s^2}.
\end{equation}

Per il segmento finale relativo al colletto vengono utilizzati valori inferiori:

\begin{equation}
    v_{\max}=0.03\,\mathrm{m/s},
    \qquad
    a_{\max}=0.002\,\mathrm{m/s^2}.
\end{equation}

Per il primo segmento di avvicinamento e per il raccordo identificato dal codice con indice sei vengono invece assegnati valori più elevati:

\begin{equation}
    v_{\max}=0.2\,\mathrm{m/s},
    \qquad
    a_{\max}=0.06\,\mathrm{m/s^2}.
\end{equation}

La riduzione dei limiti cinematici sul colletto è coerente con la scelta di percorrere più lentamente un tratto che richiede maggiore attenzione geometrica. Si tratta, tuttavia, di una scelta progettuale e non di un limite fisico universale: i valori devono essere verificati sperimentalmente rispetto alla precisione richiesta, alla risposta del robot e all'interazione con il tessuto.

\section{Calcolo della velocità cartesiana}

Una volta definita la legge temporale, è possibile determinare la velocità cartesiana desiderata dell'organo terminale.

Dalla relazione

\begin{equation}
    \mathbf{x}(t)=\mathbf{r}(u(t)),
\end{equation}

applicando la regola della derivazione composta si ottiene

\begin{equation}
    \dot{\mathbf{x}}(t)
    =\frac{d\mathbf{r}}{du}\frac{du}{dt}.
    \label{eq:velocita_catena}
\end{equation}

La velocità cartesiana è diretta lungo la tangente alla curva. Per imporre che il suo modulo coincida con la velocità scalare pianificata $v(t)$, viene utilizzato il versore tangente

\begin{equation}
    \mathbf{t}(u)=
    \frac{\mathbf{r}'(u)}
    {\|\mathbf{r}'(u)\|}.
\end{equation}

La velocità desiderata può quindi essere scritta come

\begin{equation}
    \dot{\mathbf{x}}_d(t)
    =v(t)\mathbf{t}(u(t)).
    \label{eq:velocita_cartesiana}
\end{equation}

Esplicitando le componenti, si ottiene

\begin{equation}
    v_x(t)=
    v(t)\frac{x'(u)}{\|\mathbf{r}'(u)\|},
\end{equation}

\begin{equation}
    v_y(t)=
    v(t)\frac{y'(u)}{\|\mathbf{r}'(u)\|},
\end{equation}

\begin{equation}
    v_z(t)=
    v(t)\frac{z'(u)}{\|\mathbf{r}'(u)\|}.
\end{equation}

Questa formulazione è importante dal punto di vista fisico perché separa il modulo della velocità dalla direzione del movimento. Il profilo temporale stabilisce quanto rapidamente avanzare, mentre la derivata della spline determina in quale direzione muoversi.

La norma della velocità cartesiana risulta idealmente pari a $v(t)$:

\begin{equation}
    \|\dot{\mathbf{x}}_d(t)\|=v(t),
\end{equation}

purché la derivata della curva non sia nulla.

\section{Pianificazione dell'orientamento dell'organo terminale}

Oltre alla posizione, è necessario definire l'orientamento dell'organo terminale durante il movimento. Tale orientamento determina la direzione degli assi locali dell'utensile e contribuisce a stabilire come esso si presenta rispetto alla superficie della maglia.

Per i segmenti senza un orientamento imposto, il codice costruisce una terna cartesiana a partire dalla direzione tangente al percorso. Per ogni punto campionato, la direzione tangente viene utilizzata come asse locale $x_e$:

\begin{equation}
    \mathbf{x}_e=
    \frac{\mathbf{r}'(u)}
    {\|\mathbf{r}'(u)\|}.
\end{equation}

Viene quindi introdotto un vettore di riferimento verticale

\begin{equation}
    \mathbf{z}_{\mathrm{ref}}=
    \begin{bmatrix}
        0 & 0 & -1
    \end{bmatrix}^{T}.
\end{equation}

Il secondo asse della terna viene calcolato mediante il prodotto vettoriale:

\begin{equation}
    \mathbf{y}_e=
    \frac{\mathbf{z}_{\mathrm{ref}}\times\mathbf{x}_e}
    {\|\mathbf{z}_{\mathrm{ref}}\times\mathbf{x}_e\|}.
\end{equation}

Il terzo asse viene infine ricostruito per ottenere una terna ortogonale:

\begin{equation}
    \mathbf{z}_e=\mathbf{x}_e\times\mathbf{y}_e.
\end{equation}

La matrice di orientamento desiderata è quindi

\begin{equation}
    R_d=
    \begin{bmatrix}
        \mathbf{x}_e & \mathbf{y}_e & \mathbf{z}_e
    \end{bmatrix}.
    \label{eq:matrice_orientamento}
\end{equation}

Dal punto di vista geometrico, questa costruzione allinea il primo asse locale alla direzione di avanzamento e utilizza la verticale come riferimento per definire gli altri due assi. In questo modo l'orientamento varia in relazione alla forma della curva, anziché essere mantenuto arbitrariamente costante.

Quando il prodotto vettoriale impiegato per calcolare $\mathbf{y}_e$ ha norma molto piccola, la direzione tangente è quasi parallela al vettore verticale di riferimento. In tale condizione la costruzione diventa numericamente mal condizionata; il codice prevede un vettore alternativo per evitare una divisione per una norma prossima a zero.

\subsection{Transizioni tra orientamenti}

Nei raccordi tra segmenti, l'orientamento iniziale e quello finale possono essere differenti. Per evitare una variazione improvvisa della posa, viene utilizzata l'interpolazione sferica dei quaternioni, nota come SLERP (\textit{Spherical Linear Interpolation}).

Indicando con $R_0$ e $R_1$ le matrici di orientamento agli estremi del raccordo, si convertono tali orientamenti in rappresentazioni quaternioniche e si interpola tra le due pose in funzione del parametro normalizzato $u$:

\begin{equation}
    R_d(u)=\operatorname{SLERP}(R_0,R_1,u).
\end{equation}

La SLERP segue un percorso regolare nello spazio delle rotazioni e consente di evitare l'interpolazione indipendente degli angoli di Eulero, che può produrre comportamenti poco regolari.

La transizione dell'orientamento viene quindi coordinata con l'avanzamento lungo il raccordo. La posizione segue la spline cartesiana, mentre l'orientamento evolve gradualmente tra le due configurazioni estreme.

\section{Concatenazione della traiettoria completa}

Ogni segmento produce una sequenza temporale di posizioni, velocità cartesiane e orientamenti. Al termine della pianificazione, le sequenze vengono concatenate mantenendo l'ordine dei tratti definito in precedenza.

La traiettoria risultante può essere rappresentata come

\begin{equation}
    \mathcal{T}=
    \left\{
    \mathbf{x}_d(t_i),
    \dot{\mathbf{x}}_d(t_i),
    R_d(t_i)
    \right\}_{i=0}^{N-1},
\end{equation}

dove $N$ è il numero totale di campioni.

Per ogni istante campionato sono quindi disponibili:
\begin{itemize}
    \item la posizione cartesiana desiderata $\mathbf{x}_d(t_i)$;
    \item la velocità lineare desiderata $\dot{\mathbf{x}}_d(t_i)$;
    \item l'orientamento desiderato $R_d(t_i)$.
\end{itemize}

La pianificazione produce così una descrizione cinematica completa del movimento nello spazio cartesiano. La traiettoria non specifica soltanto quali punti attraversare, ma anche come procedere lungo il percorso e come orientare l'organo terminale durante ciascun tratto.

\section{Considerazioni fisiche e limiti del modello}

La metodologia adottata consente di controllare separatamente la geometria del percorso e la sua evoluzione temporale. L'interpolazione con spline cubiche riduce le discontinuità geometriche all'interno dei singoli segmenti, mentre i profili di velocità e accelerazione limitano la rapidità con cui il percorso viene attraversato.

Questi aspetti sono particolarmente rilevanti nel contesto della manipolazione di un tessuto. Un movimento troppo rapido o caratterizzato da variazioni brusche può aumentare le sollecitazioni dinamiche e compromettere la precisione con cui l'utensile segue la linea desiderata. Al contrario, una velocità più contenuta nei tratti delicati permette di adottare una strategia di movimento più prudente.

È tuttavia importante distinguere la pianificazione cinematica dalla modellazione dinamica dell'interazione con il tessuto. Il profilo di velocità limita la velocità tangenziale e l'accelerazione lungo il percorso, ma non calcola direttamente le forze di contatto, la tensione della maglia o la sua deformazione. Inoltre, l'accelerazione cartesiana totale non coincide necessariamente con la sola accelerazione tangenziale, poiché la curvatura del percorso introduce anche una componente normale:

\begin{equation}
    \mathbf{a}(t)=
    a_t(t)\mathbf{t}(t)
    +v(t)^2\kappa(t)\mathbf{n}(t),
\end{equation}

dove $\kappa(t)$ è la curvatura locale e $\mathbf{n}(t)$ è il versore normale alla traiettoria. Di conseguenza, anche quando il modulo dell'accelerazione tangenziale è limitato, curve molto strette possono generare accelerazioni normali significative.

Infine, la pianificazione viene eseguita separatamente per ciascun segmento. La continuità della posizione tra i tratti è favorita dalla costruzione dei punti di raccordo, ma la continuità della velocità e dell'accelerazione ai confini tra segmenti dipende anche dai profili temporali assegnati e dalla loro concatenazione. Per questo motivo, la verifica della traiettoria completa mediante grafici di posizione, velocità e orientamento costituisce un passaggio importante per individuare eventuali variazioni indesiderate.

In sintesi, il metodo consente di ottenere un percorso cartesiano regolare e parametrizzato nel tempo, con una velocità adattata alle diverse fasi del movimento. La qualità fisica del movimento dipende tuttavia anche dalla curvatura locale, dalla continuità tra segmenti e dall'effettiva interazione tra organo terminale e tessuto.
